\chapter{Teoría de números}
En este capítulo hablaremos de
lo básico que concierne a la
teoría de números de olimpiadas.

No discurtiremos demasiado los
problemas que aparecen en omapa,
pero si lo elemental que se aplican
a éstos e innumerables problemas
de práctica.

\section{Divisibilidad}
La intuición de este concepto
se esclarifica con ejemplos.
Decimos que un entero positivo $a$
es divisible entre otro entero
positivo $n$ si es posible
partir el número $a$ en $n$ partes.
Por ejemplo se nos ocurre que definitivamente
$2$ divide a $4$ ya que es posible
partir el $4$ en $2$ partes iguales.
De manera similar $6$ es divisible
entre $2$ ya que es posible partir
el $6$ en $3 + 3$.
También nos damos cuenta de que
\[6 = 2 + 2 + 2\]
por lo que también es posible partir
al $6$ en $3$ partes iguales, o sea
$3$ divide a $6$.

Por lo que podemos considerar
la siguiente manera de pensar en la
división:

\begin{intuition*}
	Un entero positivo $a$ divide a
	un entero positivo $b$ si es posible
	partir a $b$ en $a$ partes iguales.
\end{intuition*}

Este concepto se puede generalizar
con la misma idea de la siguiente manera:

\begin{definition}
	\label{definition:div}
	Decimos que un entero $a$ divide
	a otro entero $b$ si es posible
	encontrar un entero $z$ tal que
	\[az = b.\]
\end{definition}

\begin{question*}
	¿Cómo se relaciona esto con que
	se posible partir un entero en
	una cantidad entera de partes?
\end{question*}

Notemos también que el número $1$
divide a todos los enteros, porque
es obvio que es posible partir un número
en una sola parte.

A la relación de divisibilidad
se la denota por el símbolo `$\mid$'
y ``$a$ divide a $b$'' se simboliza como
\[a \mid b.\]

De manera similar si $a$ \textbf{no}
divide a $b$ se simboliza como
\[a \nmid b.\]

\subsection{Propiedades de la división}
Ahora establecemos algunas propiedades
que son fundamentales relacionados a
la división entre dos números enteros.

\begin{proposition}[Transitividad de la divisibilidad]
	\label{proposition:transitivitydiv}
	Si $a \mid b$ y $b \mid c$, entonces $a \mid c$
\end{proposition}

\begin{proposition}[Desigualdad con la divisibilidad]
	Si $a \mid b$, entonces $\abs a \leq \abs b$.
\end{proposition}

Esto tiene sentido ya que
claramente no es posible partir un
entero en más partes que unidades éste
contiene.

\begin{proposition}
	Si $a \mid b$ y $a \mid b + c$, entonces
	$a \mid c$.
\end{proposition}

\begin{proposition}
	Si $a \mid b$, entonces $a \mid zb$
	para todos los enteros $z$.
\end{proposition}

\begin{proposition}
	Si $a \mid b$ si y solo si $a \mid -b$
	si y solo si $-a \mid b$.
\end{proposition}

\subsection*{Problemas para esta sección}
Estos problemas se hacen todos
con la \autoref{definition:div}.

\begin{problem}
Demuestre las últimas cinco
proposiciones.
\end{problem}

\section{Algoritmo de división Euclidiana}
Este el el algoritmo básico en el que se sostiene
toda la teoría de números.

Este es un teorema que nos permite expresar
cada par de dividendo y divisor $a$ y $b$
como
\[a = bq + r\]
donde $0 \leq r < b$.
Que en síntesis

\section{Números primos}
Hay enteros positivos que no son el
producto de otros dos enteros mayores que $1$,
como por ejemplo el $2$ y el $3$,
vemos que no existen otros enteros
distintos de $1$ que \textit{dividen}
a $2$ a parte de sí mismo.

Esto también pasa con el $5$, revisando número
por número vemos que ninguno divide a este
salvo el $1$ y el $5$. De hecho, hay infinitos
números que cumplen esto, y se los
llama \textit{números primos}.

\begin{definition}
	Un número primo es un entero positivo con
	exactamente dos divisores positivos.
\end{definition}

Como ya vimos, el número $1$ divide a todo número,
y a su vez $n$ siempre divide a $n$, esto implica
que si $n$ es primo, entonces éstos son sus únicos
dos divisores positivos (nótese que así el $1$ no es primo).

\begin{theorem}[Teorema fundamental de la aritmética]
	Todo entero positivo $n$ puede ser descompuesto
	en factores primos de manera única.
\end{theorem}

Para desglosar qué dice este teorema,
una manera de verlo es que si $n$ es un
entero positivo y $p_1$, $p_2$, \dots, $p_n$
son todos los primos que lo dividen, entonces
existen unos únicos exponenetes $\alpha_1$, $\alpha_2$,
\dots, $\alpha_n$ tales que
\[n = p_1^{\alpha_1}p_2^{\alpha_2}\cdots p_n^{\alpha_n}.\]

Se puede afirmar con bastante seguridad que
este facto se aplica a casi todos los problemas
de teoría de números.

\section{Aritmética modular}
Esta es una idea en la que se sostiene
casi todo lo desarrollado respecto
a la tería de números, y consiste en
categorizar qué en general
si analizamos los restos que deja
el algoritmo de división euclidiana.

Para nosotros este pedazo de teoría
nos ayudará a abreviar muchas o

El clásico ejemplo para describir a la
aritmética modular es el reloj de $12$ horas:

% TODO: Insertar imagen de reloj de 12 horas

Digamos que el reloj empieza en la hora $0$,
entonces al cabo de $10$ horas el reloj
marcará, $10$, al cabo de $11$ marcará $11$,
pero al paso de $12$ horas éste volverá a marcar $0$.
\footnote{Esquivada olímpica.} % TODO: Quitar este footnote.

Esto se puede formalizar a través de
analizar los restos en \textit{módulo} $12$
y decimos que las horas que marca el reloj
son los enteros \textit{módulo} $12$.

\begin{definition}
	Dos enteros son congruentes módulo $n$
	si tienen el mismo resto divididos entre $n$.
\end{definition}

Esto abre paso a, quizás, querer separar
a los enteros teniendo en cuenta
su resto al dividirse entre $n$,
en particular, no es difícil de ver que:

\begin{theorem}\label{theorem:congruence}
	Dos enteros $a$ y $b$ son congruentes
	módulo $n$ si y solo si $n$ divide a $a - b$.
\end{theorem}

Más aún, ver a los enteros desde esta perspectiva
nos abre un gran abanico de herramientas
para entenderlos, entre ellos se dejan ciertas
propiedades como ejercicios al lector.

\begin{definition}
	Denotamos a `$\equiv$' como el símbolo
	de congruencia. La expresión
	``$a$ y $b$ son congruentes módulo $n$''
	se escribe matemáticamente de la siguiente
	manera:
	\[a \equiv b \pmod n\]
\end{definition}

\subsection*{Problemas para esta sección. }
\begin{problem}
Demuestre el \autoref{theorem:congruence}.
\end{problem}

\begin{problem}[Transitividad de la congruencia]
Demuestre que si $a \equiv b$ y $b \equiv c$,
entonces $a \equiv c$.
\end{problem}

\begin{problem}[Suma de congruencias]
Demuestre que si $a \equiv b$ y $x \equiv y$ en
módulo $n$, entonces
\[a + x \equiv b + y \pmod n.\]
\end{problem}

\begin{problem}[Congruencia y multiplicaciones]
\label{theorem:congprod}
Demuestre que si $a \equiv x$ y $b \equiv y$
en módulo $n$ entonces
\[ab \equiv xy \pmod n.\]
\end{problem}

\begin{problem}[Congruencia y la exponenciación]
Demuestre que si $a \equiv b$ en módulo $n$,
entonces
\[a^k \equiv b^k \pmod n\]
para todo entero positivo $k$.

\begin{hints}
	\begin{hint}
		\autoref{theorem:congprod}.
	\end{hint}
\end{hints}
\end{problem}

\section{Problemas}
\begin{problem}
Alicia y Beto juegan a un juego
por turnos de forma alternada.
En cada turno el jugador
escribe un dígito del $1$
al $9$ en la pizarra.
Si ambos juegan la misma cantidad de veces
y empieza Alicia,
¿puede ella asegurar que después
de escribir todos los dígitos
sea posible escribir un número primo
con esos dígitos?

\begin{hints}
	\begin{hint}
		¿Puede Beto asegurar que el número sea
		múltiplo de algo?
	\end{hint}

	\begin{hint}
		Múltiplos de $3$.
	\end{hint}
\end{hints}
\end{problem}

\begin{problem}
Calcule la cantidad de divisores enteros positivos de $2026$.
\end{problem}

\begin{problem}
Calcule la suma de los divisores naturales de $2026$.
\end{problem}

\begin{problem}
Encuentre la cantidad de triplas de enteros positivos $a$, $b$ y $c$ tales que $a$ divide a $bc + 1$, $b$ divide
a $ac + 2$ y $c$ divide a $ab + 3$.
\end{problem}

\begin{problem}
Encuentre el resto de dividir $2026^{2026^{2026}}$ entre $7$.
\end{problem}

\begin{problem}
Encuentre todas las ternas $a$, $b$, $c$ de enteros positivos
tales que
\[a^2 + b^2 + 1 = c!.\]
\end{problem}

\begin{problem}
Demuestre que que para todo $n$ natural se cumple que
\[6 \mid n^3 + 5n.\]
\end{problem}

\begin{problem}
Encuentre todos los pares de enteros $\left(x, y\right)$ tales que
\[\frac{xy}{x + y} = 144\]
\end{problem}

\begin{problem}[CAM JT 2025]
Demuestre que si $\overline{abc\dots{x}yz}$ es múltiplo de $37$,
entonces $\overline{z0y0x0\dots0c0b0a}$ es múltiplo de 37.
(Ej.: $259$ es múltiplo de $37$, entonces $90502$ es múltiplo de 37).
\end{problem}

\begin{problem}
Encuentre todas las funciones $f\colon\NN\to\NN$ tales que $f(2) = 2$, $f$
es estrictamente creciente y $f(xy) = f(x)f(y)$ para todos los pares de
enteros positivos $(x, y)$.
\end{problem}
